\documentclass[11pt,a4paper]{article}
\usepackage[utf8]{inputenc}
\usepackage[ngerman]{babel}
\usepackage{amsmath,amssymb,amsthm}
\usepackage{geometry}
\usepackage{booktabs}
\usepackage{xcolor}
\usepackage{tcolorbox}

% Seitenränder definieren
\geometry{left=2.5cm, right=2.5cm, top=2.5cm, bottom=2.5cm}

% Farbdefinitionen für Hervorhebungen
\definecolor{blockbg}{HTML}{F8F9FA}
\definecolor{blockborder}{HTML}{D1D5DB}
\definecolor{stuckbg}{HTML}{FEF2F2}
\definecolor{stuckborder}{HTML}{EF4444}
\definecolor{stucktext}{HTML}{991B1B}

\tcbuselibrary{skins,breakable}

% Box für Givens/Goal Tabellen
\newtcolorbox{vellemantable}[1][]{
  colback=blockbg,
  colframe=blockborder,
  arc=3mm,
  boxrule=1pt,
  left=2mm, right=2mm, top=2mm, bottom=2mm,
  #1
}

% Box für die Stellen, an denen der Beweis scheitert
\newtcolorbox{stuckbox}[1][]{
  colback=stuckbg,
  colframe=stuckborder,
  coltitle=stucktext,
  fonttitle=\bfseries,
  arc=3mm,
  boxrule=1.2pt,
  title={Hier bleiben wir stecken!},
  #1
}

\title{\textbf{Untersuchung von Bild- und Urbildmengen}\\[0.3em]
\large Formalismen nach Daniel Velleman (\textit{How to Prove It})}
\author{}
\date{}

\begin{document}

\maketitle

\section*{Problemstellung}
Sei $f: X \to Y$ eine beliebige Funktion, $A \subseteq X$ und $B \subseteq Y$. Wir untersuchen, warum die folgenden Gleichheiten im Allgemeinen \textbf{nicht} gelten und an welchen exakten Schritten der Beweis ohne Zusatzvoraussetzungen scheitert.

---

\section{Analyse von $f(f^{-1}(B)) = B$}

Allgemein gilt nur $f(f^{-1}(B)) \subseteq B$. Der Versuch, die Rückrichtung $B \subseteq f(f^{-1}(B))$ zu beweisen, verläuft wie folgt:

\subsubsection*{Ziel}
Zeige $\forall y \bigl(y \in B \implies y \in f(f^{-1}(B))\bigr)$.

\begin{vellemantable}
\centering
\begin{tabular}{cll}
\textbf{Schritt} & \textbf{Givens} & \textbf{Goal} \\ \hline
1 & $f: X \to Y, \; B \subseteq Y$ & $B \subseteq f(f^{-1}(B))$ \\
2 & $y \in B$ & $y \in f(f^{-1}(B))$ \\
3 & $y \in B$ & $\exists x \in f^{-1}(B) \: (f(x) = y)$ \quad \textit{(Def. Bild)} \\
4 & $y \in B$ & $\exists x \: (f(x) \in B \land f(x) = y)$ \quad \textit{(Def. Urbild)} \\
\end{tabular}
\end{vellemantable}

\begin{stuckbox}
Aus $y \in B$ folgt \textbf{nicht}, dass überhaupt ein $x \in X$ existiert mit $f(x) = y$.

Wenn $y$ nicht im Bild von $f$ liegt ($\text{Bild}(f)$), existiert kein solches $x$. Das Urbild $f^{-1}(B)$ enthält für dieses $y$ kein Element, und $y$ kann nicht erneut von $f$ getroffen werden.

\medskip
\textbf{Benötigte Eigenschaft:} \textbf{Surjektivität} von $f$ (bzw. $B \subseteq \text{Bild}(f)$), um die Existenz eines $x \in X$ mit $f(x) = y$ für jedes $y \in B$ zu garantieren.
\end{stuckbox}

---

\section{Analyse von $f^{-1}(f(A)) = A$}

Allgemein gilt nur $A \subseteq f^{-1}(f(A))$. Der Versuch, die Rückrichtung $f^{-1}(f(A)) \subseteq A$ zu beweisen, verläuft wie folgt:

\subsubsection*{Ziel}
Zeige $\forall x \bigl(x \in f^{-1}(f(A)) \implies x \in A\bigr)$.

\begin{vellemantable}
\centering
\begin{tabular}{cll}
\textbf{Schritt} & \textbf{Givens} & \textbf{Goal} \\ \hline
1 & $f: X \to Y, \; A \subseteq X$ & $f^{-1}(f(A)) \subseteq A$ \\
2 & $x \in f^{-1}(f(A))$ & $x \in A$ \\
3 & $f(x) \in f(A)$ \quad \textit{(Def. Urbild)} & $x \in A$ \\
4 & $\exists a \in A \: (f(a) = f(x))$ \quad \textit{(Def. Bild)} & $x \in A$ \\
5 & $a_0 \in A \land f(a_0) = f(x)$ & $x \in A$ \\
\end{tabular}
\end{vellemantable}

\begin{stuckbox}
Wir wissen nur, dass $x$ denselben Funktionswert hat wie irgendein Element $a_0 \in A$ (also $f(x) = f(a_0)$).

Das bedeutet jedoch nicht automatisch, dass $x = a_0$ gilt. Es könnte ein weiteres Element $x \notin A$ geben, das auf denselben Bildpunkt abgebildet wird.

\medskip
\textbf{Benötigte Eigenschaft:} \textbf{Injektivität} von $f$, da hieraus direkt $f(a_0) = f(x) \implies a_0 = x$ folgt (womit $x = a_0 \in A$ bewiesen wäre).
\end{stuckbox}

\end{document}
